\scanpage{12}
\uline{Corollaire} 2.14. $\underline C$ étant comme ci-dessus la catégorie des $A$-modules de type fini annulés par une puissance de $I$, soit $T:\underline C^o\handwritten{\to}(\mathrm{Ab})$ un foncteur
contravariant exact à gauche de $\underline C$ dans $(\mathrm{Ab})$, Soit $(Y_i)_{i\handwritten{\in}I}$ la famille des
composantes irréductibles de $Y\inserted{=V(I)}$, et soit $(Y_i)_{i\handwritten{\in}J}$ une sous-famille telle
que pour tout $M\handwritten{\in}\mathrm{Ob}\,\underline C$, la nullité de $T(M)$ (qui de toute façon ne dépend
que de $\operatorname{supp}M$ en vertu de 2.4.) soit équivalente à la relation ``$Y_i\handwritten{\not\subset}
\operatorname{supp}M$ pour tout $i\handwritten{\in}J$''. Soit $S\inserted{\handwritten{=(y_i)_{i\in J}}}$ l'ensemble des points génériques des
$Y_i$, $i\handwritten{\in}J$. Alors pour tout $M\handwritten{\in}\mathrm{Ob}\,\underline C$, on a
\[
\operatorname{Ass}T(M)=\operatorname{supp}M\handwritten{\cap} S\quad,
\]
en particulier le $A$-module $T(M)$ est sans idéaux premiers associés imm.

En effet, posant $H=\lim T(A/I^n)$, on a un isomorphisme
\[
T(M)\handwritten{\simeq}\operatorname{Hom}(M,H)
\]
d'où
\[
\operatorname{Ass}\deleted{\unsure{E|P}}T(M)=\operatorname{supp}M\handwritten{\cap}\operatorname{Ass}H\handwritten{,}\tag*{\%}
\]
et il reste à prouver la formule
\[
\operatorname{Ass}H=S\quad.
\]
\deleted{Soit donc} Evidemment $\operatorname{Ass}H\handwritten{\subset}\operatorname{supp}H\handwritten{\subset}Y$, il reste à prouver que si $y\handwritten{\in}Y$,
alors $y\handwritten{\in}\operatorname{Ass}H$ équivaut à $y\handwritten{\in}S$. Or soit $M=A/\underline p$, où $\underline p=\underline p_y$, \inserted{\handwritten{donc $V(M)=V(\underline p)=\bar y$,}} en
vertu de \% on a $T(M)\handwritten{\ne}0$ (i.e. $\operatorname{Ass}T(M)\handwritten{\ne\varnothing}$\handwritten{)} sss $\deleted{\unsure{V(\underline p)}}\handwritten{\bar y}$ \deleted{\unsure{$=\bar y$}} renoontre \deleted{pas} $\operatorname{Ass}H$,
or par hypothèse on a $T(M)\handwritten{\ne}0$ sss \deleted{$\operatorname{supp}M=V(\underline p)$}\inserted{\handwritten{$\bar y$}} \deleted{\unsure{ne}} rencontre \deleted{pas} $S$,
ou ce qui revient au même, sss $\deleted{\unsure{\underline p}}\handwritten{y\in}S$. Cela prouve d'abord que $\operatorname{Ass}H\handwritten{\subset}S$,
(prendre $y\handwritten{\in}\operatorname{Ass}H$), puis $S\handwritten{\subset}\operatorname{Ass}H$ puisque tout $y\handwritten{\in}S$ doit avoir une
spécialisation dans $S$, donc être dans $S$ puisque deux éléments de $S$ dont
\deleted{sans} l'un spécialise l'autre sont identiques.

Le résultat 2.10. nous incite à traduire les propriétés des
modules \inserted{$H$} de support $Y=V(I)$ en termes du foncteur exact à gauche
correspondant $M\handwritten{\leadsto}\operatorname{Hom}(M,\mathbb H)$ sur $\underline C$. Ici, nous nous contentons d'expliciter
le résultat suivant:

\uline{Proposition} 2.15. Avec les notations précédentes, pour que $H$ soit un
$A$-module injectif, il f et s que le foncteur $T_H(M)=\deleted{\unsure{H}}\operatorname{Hom}(M,H)$ sur la
catégorie $\underline C$ des $A$-modules de type fini \deleted{à supp\unsure{o}} annulés par une puissance
de $I$, soit exact à droite donc exact. En d'autres termes, 2.10. induit
une équivalence entre la catégorie des foncteurs contravariants \uline{exacts}
$T$ de $C$ dans$(\mathrm{Ab})$, avec la catégorie des $A$-modules injectifs à support $\handwritten{\subset}V(I)$
i.e. dont tout élément est annulé par une puissance de $I$.
% Faint inverted offset from the neighbouring typescript is visible beneath the final paragraph.
% archive sheet number: 11

